\section{Conclusion and Future Work}\label{sec:conclusion}
This study has established that the reliability of few-shot anomaly thresholds depends on both the amount and the statistical unit of normal evidence. At nominal $\alpha=0.20$, target-only LOIO reaches an empirical FAR of 0.341 on Gaussian-corrupted MVTec at $k=4$, showing that an attainable operating point need not remain stable under shift. The finite floor $1/(k+1)$ explains when a target-only rule is structurally silent, but crossing that floor does not supply an exchangeability guarantee.

The category-count analysis has quantified the separate obstacle faced by source-assisted certification. With all category losses zero, optimistic multiplicity-free certification by a deterministic, uniformly valid, distribution-free 95\% UCB requires at least 14, 29, and 59 iid category draws at $\alpha=0.20$, 0.10, and 0.05. These are necessary lower limits, not sufficient budgets for CRESS\@. The frozen allocation and Hoeffding rule require more, whereas the strict audit provides only three or four certification categories. Even a certification-heavy reallocation with nonempty reference and proposal roles provides at most 9 to 13 categories, so changing the split cannot meet the optimistic 14-category limit. Category-unit certification consequently assigns $\tau^\star=0$ to every target cell across all 960 frozen configurations, which share the same boundary and are not independent trials. In contrast, image-unit analyses under the idealized selected-mixture iid model select positive thresholds in 36.7\% to 60.3\% of target cells, but target the selected source mixture rather than marginal risk over a new-category draw.

The practical contribution is therefore a study-design discipline: distinguish ranking from threshold reliability, report attainable rather than nominal operating points, audit normal false alarms under declared shifts, define the certification unit before observing results, and justify iid category sampling when the claim concerns a new category. Future work should collect category-rich source archives or introduce explicit parametric, hierarchical, or side-information assumptions with independently verified validity. It should also test real production shifts, temporal dependence, condition-estimation error, and end-to-end deployment cost. These directions can improve usefulness, but they should not relabel repeated images as category evidence, confuse marginal with category-conditional risk, or convert a source-domain certificate into unconditional target control.
